A deployed classifier needs different responses to covariate, label and concept shift: keep the model, correct its class prior, or retrain it. We study sequential monitoring that reports which response is warranted, with error control, when only a small random fraction of labels is revealed. Typed Shift Sequential Diagnosis (TSSD) runs three witnesses on one anytime-valid budget: a paired kernel witness for inputs, a gated prevalence witness, and a new posterior-validity witness. The last is a universal-inference e-process whose composite null contains every posterior reachable by a prior and tempering correction of the frozen model, under arbitrary covariate shift and with a label shift at an unknown time. Its crossing is therefore a certified claim that prior correction is insufficient. Every witness also runs as a Shiryaev--Roberts e-detector with average-run-length and finite-window false-alarm bounds. We compare TSSD on identical streams with ten re-implemented detectors and a Bonferroni ensemble of specialists. On controlled streams, TSSD detects all six pure and mixed shift types at a 1.3\% null alarm rate and types 180--199 of 200 streams correctly per type. The ensemble types 0 of 194 covariate shifts correctly and makes 149 false concept claims, where TSSD makes none in any experiment. On chronological electricity data, TSSD certifies concept drift in all 260 runs on segments where only a refit repairs the model, and in none where a prior correction suffices. TSSD is slower than loss-based specialists on concept alternatives, and its certified claim needs labels.
